\documentclass[runningheads]{llncs}

\usepackage[T1]{fontenc}
\usepackage{graphicx}
\usepackage{amsmath}
\usepackage{booktabs}
\usepackage{url}

\begin{document}

\title{Fine-Tuning YOLO11n-OBB for Oriented Vehicle Detection in Peruvian Traffic Scenes}
\titlerunning{YOLO11n-OBB for Peruvian Traffic Scenes}

% Replace these fields with the final author information before camera-ready submission.
\author{Anonymous Author(s)}
\authorrunning{Anonymous Author(s)}
\institute{Anonymous Institution, Peru\\
\email{anonymous@example.org}}

\maketitle

\begin{abstract}
Reliable traffic monitoring in resource-constrained urban settings requires
detectors that remain accurate under oblique viewpoints while being suitable for
deployment on modest hardware. This work fine-tunes YOLO11n-OBB for nine vehicle
categories in Peruvian traffic imagery. We build clip-level train, validation,
and test partitions to prevent adjacent video frames from appearing in different
splits. A controlled data-scaling study evaluates training subsets of 1,000,
1,500, 2,000, and 2,500 images using the same 50-epoch protocol and an effective
input size of 544 pixels. Test mAP@50 increases from 66.2\% with 1,000 images to
91.4\% with 2,500 images, while mAP@50--95 rises from 54.8\% to 78.5\%.
On the 2,500-image subset, a second configuration trained for 70 epochs at 640
pixels achieves 93.9\% mAP@50 and 81.7\% mAP@50--95 with 2.66 million parameters.
Desktop CPU measurements are reported only as a four-thread proxy; the NCNN
export is prepared for a separate on-device Raspberry Pi evaluation. The results
show that a compact oriented detector can obtain strong accuracy from a modest
domain-specific dataset, and establish a reproducible basis for low-cost traffic
analytics in Peru.

\keywords{Oriented object detection \and Vehicle detection \and YOLO11n-OBB
\and Edge AI \and Intelligent transportation systems.}
\end{abstract}

\section{Introduction}

Traffic monitoring supports vehicle counting, incident analysis, and urban
planning. Yet the deployment conditions of many cities make it impractical to
assume dedicated GPUs at every camera. A useful detector must therefore balance
accuracy with a small computational footprint. This is particularly relevant in
Peruvian road scenes, where surveillance footage can contain cars, buses,
microbuses, mototaxis, motorcycles, and trucks under substantial perspective
distortion.

Conventional axis-aligned boxes can contain considerable background when a
vehicle is viewed at an angle. In dense traffic, they can also overlap heavily
with neighboring vehicles. Oriented bounding boxes (OBBs) encode the extent and
orientation of an object explicitly, an advantage already exploited in aerial
image benchmarks such as DOTA~\cite{xia2018dota,ding2022dota}. The same geometry
is relevant to elevated traffic cameras, although the visual domain, classes,
and practical deployment constraints differ from aerial imagery.

This paper investigates the fine-tuning of YOLO11n-OBB, a compact member of the
Ultralytics YOLO family that supports OBB prediction~\cite{ultralytics_yolo11,ultralytics_obb}. We focus on the amount of domain-specific data required for
useful performance and on a deployment path that can later be assessed on ARM
hardware. The present paper makes the following contributions:

\begin{itemize}
\item A reproducible conversion and clip-level splitting protocol for nine
vehicle categories from Peruvian traffic videos.
\item A controlled data-scaling experiment from 1,000 to 2,500 images under a
fixed 50-epoch training protocol.
\item A higher-resolution 2,500-image configuration that attains 93.9\%
mAP@50 and 81.7\% mAP@50--95 with 2.66 million parameters.
\item A pre-specified NCNN benchmark protocol for Raspberry Pi deployment that
keeps desktop CPU measurements separate from device-side results.
\end{itemize}

\section{Related Work}

\subsection{Vehicle Detection and Oriented Bounding Boxes}

One-stage detectors popularized the use of direct, real-time object prediction
from images~\cite{redmon2016yolo}. Vehicle benchmarks such as UA-DETRAC have
made it possible to evaluate detection and tracking in traffic video, but they
do not address the class composition and OBB annotations used in this
work~\cite{wen2015uadetrac}. In aerial imagery, arbitrary-oriented quadrilaterals
are essential because objects exhibit strong variation in rotation and scale;
DOTA provides a widely used benchmark for this setting~\cite{xia2018dota}.

The current study transfers this representation to elevated traffic footage.
OBBs are not assumed to solve every source of error---occlusion, small objects,
and class imbalance remain difficult---but they provide a closer geometric
description of rotated vehicles than horizontal boxes. YOLO11 supports OBB
training and inference in the Ultralytics implementation used here
\cite{ultralytics_obb}.

\subsection{Efficient Inference for Edge Devices}

Embedded deployment requires that model conversion and runtime behavior be
reported separately from training accuracy. NCNN is an inference framework aimed
at mobile, embedded, and desktop environments, and provides a conversion path
from PyTorch or ONNX models~\cite{ncnn}. We export the selected checkpoint to
NCNN, but do not claim Raspberry Pi throughput until it has been measured on the
target device. This distinction is important because a desktop CPU measurement,
even with a restricted number of threads, cannot characterize an ARM platform.

\section{Dataset and Experimental Method}

\subsection{Dataset and OBB Conversion}

The data consist of traffic surveillance images from Peru with annotations for
nine categories: \textit{car}, \textit{van}, \textit{microbus}, \textit{minibus},
\textit{bus}, \textit{articulated\_bus}, \textit{truck}, \textit{mototaxi}, and
\textit{motorcycle}. Source annotations encode a class identifier, center
coordinates, width, height, and angle. Each annotation is converted to the four
normalized corner coordinates expected by YOLO-OBB.

For a box centered at $(c_x,c_y)$ with width $w$, height $h$, and orientation
$\theta$, a local corner $(d_x,d_y) \in \{(-w/2,-h/2),(w/2,-h/2),
(w/2,h/2),(-w/2,h/2)\}$ is transformed as
\begin{equation}
\begin{aligned}
x &= c_x+d_x\cos\theta-d_y\sin\theta,\\
y &= c_y+d_x\sin\theta+d_y\cos\theta.
\end{aligned}
\label{eq:obb}
\end{equation}
The resulting coordinates are divided by image width and height and clipped to
$[0,1]$.

Subsets are formed by allocating a quota to each source video and selecting
temporally separated frames within that video. Images are subsequently grouped
by clip identifier. Whole clips, rather than individual frames, are assigned to
train, validation, or test. This prevents near-duplicate adjacent frames from
leaking across partitions. The random seed for a subset of size $n$ is $42+n$.
Table~\ref{tab:splits} reports the resulting partitions.

\begin{table}
\caption{Clip-level data partitions used in the experiments.}
\label{tab:splits}
\centering
\begin{tabular}{rrrrr}
\toprule
Subset & Train & Validation & Test & Total \\
\midrule
1,000 & 700 & 150 & 150 & 1,000 \\
1,500 & 1,050 & 225 & 225 & 1,500 \\
2,000 & 1,400 & 300 & 300 & 2,000 \\
2,500 & 1,748 & 376 & 376 & 2,500 \\
\bottomrule
\end{tabular}
\end{table}

\subsection{Training Protocol and Metrics}

All experiments initialize from \texttt{yolo11n-obb.pt}. The model contains
2,655,478 parameters. The fixed scaling protocol uses a batch size of 32,
50 epochs, and a requested image size of 540 pixels. The training runtime aligns
this requested size to the 544-pixel stride-compatible effective input size;
we use the latter when naming the experimental configuration. The model with
2,500 images is also trained for 70 epochs at 640 pixels. In both cases,
patience is 20 epochs and the seed is deterministic.

The checkpoint selected by validation is evaluated once on the held-out test
split. We report precision, recall, mAP at IoU 0.50 (mAP@50), and mAP averaged
over IoU thresholds from 0.50 to 0.95 (mAP@50--95), following the common COCO
evaluation convention~\cite{lin2014coco}. CPU latency is measured on 30 test
images with batch size one and four CPU threads. It is explicitly labeled a
desktop proxy, not an edge-device result.

\section{Results}

\subsection{Data Scaling Under a Fixed Training Protocol}

Table~\ref{tab:scaling} presents the controlled data-scaling results. The
strongest result in this family is obtained with 2,500 images: 91.4\% mAP@50
and 78.5\% mAP@50--95. Relative to 1,000 images, this is a gain of 25.2 and
23.6 percentage points, respectively. The gain is not strictly monotonic at
every intermediate point: mAP@50 is slightly lower at 2,000 than at 1,500
images. Because each subset has a clip-level test split of its own, this local
variation may reflect both training-data composition and finite test-set
variation. It should not be over-interpreted.

\begin{table*}
\caption{Test performance for the controlled scaling experiment. All rows use
YOLO11n-OBB, 50 epochs, batch size 32, and a 544-pixel effective input size.
CPU measurements are a four-thread desktop proxy over 30 images.}
\label{tab:scaling}
\centering
\small
\begin{tabular}{rrrrrrr}
\toprule
Images & mAP@50 & mAP@50--95 & Precision & Recall & Latency (ms) & CPU FPS \\
\midrule
1,000 & 0.6622 & 0.5483 & 0.7742 & 0.6425 & 30.99 & 32.27 \\
1,500 & 0.8167 & 0.6769 & 0.8352 & 0.7058 & 32.11 & 31.14 \\
2,000 & 0.8129 & 0.6873 & 0.8616 & 0.7219 & 32.98 & 30.32 \\
2,500 & \textbf{0.9140} & \textbf{0.7847} & \textbf{0.9053} &
\textbf{0.8517} & 32.89 & 30.41 \\
\bottomrule
\end{tabular}
\end{table*}

\subsection{Candidate Configuration at 640 Pixels}

The 70-epoch, 640-pixel configuration is a candidate for edge deployment. It
uses the same 2,500-image subset as the last row of Table~\ref{tab:scaling},
but changes both the number of epochs and the resolution. It therefore compares
complete configurations rather than isolating either factor. As shown in
Table~\ref{tab:final}, it improves mAP@50 by 2.5 percentage points and
mAP@50--95 by 3.3 percentage points over the 50-epoch configuration.

\begin{table}
\caption{Two configurations trained on the 2,500-image subset.}
\label{tab:final}
\centering
\small
\begin{tabular}{lrrrr}
\toprule
Configuration & mAP@50 & mAP@50--95 & Precision & Recall \\
\midrule
50 epochs, 544 px & 0.9140 & 0.7847 & 0.9053 & 0.8517 \\
70 epochs, 640 px & \textbf{0.9394} & \textbf{0.8175} & \textbf{0.9281} & \textbf{0.8595} \\
\bottomrule
\end{tabular}
\end{table}

The AP@50 values of the selected 640-pixel model are shown in
Table~\ref{tab:classes}. Most vehicle categories exceed 0.90 AP@50. Motorcycle
is the lowest-scoring class (0.749), making it a useful target for future error
analysis. The present measurements do not establish the cause; likely factors
such as scale, occlusion, and class frequency should be tested rather than
assumed.

\begin{table}
\caption{Per-class AP@50 for the 70-epoch, 640-pixel model.}
\label{tab:classes}
\centering
\small
\begin{tabular}{lr@{\qquad}lr}
\toprule
Class & AP@50 & Class & AP@50 \\
\midrule
car & 0.993 & van & 0.962 \\
microbus & 0.924 & minibus & 0.941 \\
bus & 0.995 & articulated\_bus & 0.995 \\
truck & 0.981 & mototaxi & 0.915 \\
motorcycle & 0.749 & & \\
\bottomrule
\end{tabular}
\end{table}

\section{Raspberry Pi Evaluation Protocol}

The selected model has been exported to NCNN, but on-device measurements are
not yet available. To keep the deployment claim reproducible, the final version
will report the Raspberry Pi model and memory, operating system and kernel,
NCNN version, CPU thread count, thermal conditions, warm-up procedure, and the
number of timed images. Mean latency, P50, P95, and FPS will be measured for
the resolutions 640, 544, 416, and 320. The paper will also state whether the
time includes preprocessing, postprocessing, and output serialization.

Where feasible, accuracy at each input resolution will be evaluated on the same
376-image test split used for the 2,500-image model. If only latency is
measured, accuracy cells will be marked ``not evaluated'' rather than estimated.
Table~\ref{tab:edge} is intentionally included as a reporting template; it is
not evidence of Raspberry Pi performance.

\begin{table*}
\caption{Edge-benchmark reporting template. Only the desktop proxy row has
been measured.}
\label{tab:edge}
\centering
\small
\resizebox{\textwidth}{!}{%
\begin{tabular}{lrrrrrrrr}
\toprule
Runtime & Resolution & Threads & Images & Mean (ms) & P50/P95 (ms) & FPS & mAP@50 & Status \\
\midrule
PyTorch desktop CPU proxy & 640 & 4 & 30 & 33.78 & -- & 29.60 & 0.9394 & Measured \\
NCNN on Raspberry Pi & 640 & TBD & TBD & TBD & TBD & TBD & TBD & Pending \\
NCNN on Raspberry Pi & 544 & TBD & TBD & TBD & TBD & TBD & TBD & Pending \\
NCNN on Raspberry Pi & 416 & TBD & TBD & TBD & TBD & TBD & TBD & Pending \\
NCNN on Raspberry Pi & 320 & TBD & TBD & TBD & TBD & TBD & TBD & Pending \\
\bottomrule
\end{tabular}
}
\end{table*}

\section{Discussion and Limitations}

The fixed-protocol results indicate that 2,500 images are sufficient to obtain
strong held-out accuracy for the present domain. However, the scaling experiment
does not estimate uncertainty across random partitions: each data size has a
different clip-level test set. A stronger follow-up study would use a common
fixed test set, repeat the partitioning over multiple seeds, or both. It would
also report class frequencies and confidence intervals.

The study currently has no device-side result, and the desktop proxy must not
be used to claim real-time Raspberry Pi inference. In addition, the breadth of
geographic coverage, nighttime scenes, adverse weather, and camera viewpoints
must be documented before making broader generalization claims. Future work can
add these conditions and combine detection with multi-object tracking for
vehicle counting and trajectory analysis.

\section{Conclusion}

We presented a clip-level evaluation of YOLO11n-OBB for oriented vehicle
detection in Peruvian traffic scenes. Under a fixed 50-epoch protocol, increasing
the subset from 1,000 to 2,500 images improved test mAP@50 from 66.2\% to
91.4\%. A 70-epoch, 640-pixel configuration reached 93.9\% mAP@50 and 81.7\%
mAP@50--95 with 2.66 million parameters. These findings support the use of a
compact oriented detector for this domain, while the final conclusion on edge
viability awaits the planned, fully documented NCNN benchmark on Raspberry Pi.

\begin{credits}
\subsubsection{Acknowledgments} To be completed after confirming the data
license, institutional affiliations, and any funding source.

\subsubsection{Disclosure of Interests} The authors declare no competing
interests.
\end{credits}

\bibliographystyle{splncs04}
\bibliography{references}

\end{document}
